\chapter{New Frontiers}

The UT Computation Center’s history reflected a pioneering and ambitious spirit during the sixties and seventies, with a mindset of pushing boundaries, challenging the status quo, and expanding across several major initiatives. This manifested itself in concrete, systemic transformations driven by capital investments, unique administrative policies, and the direct empowerment of students and staff. A collaborative mandate to share resources regionally established the Computation Center’s role in the early digital revolution. It operated with a distinct sense of scale and regional duty. Charles Warlick, an old associate of David Young and the director of the Computation Center, encapsulated this spirit when discussing the university's large capital investments and ambitious goals for academic computing. He declared, "Texans do not undertake things in a small way. We expect large successes, and we expect to share them to the general benefit of higher education". Another of Warlick’s mottos, prominently displayed in his office, was “There they go. I had better run and catch them, for I am their leader.” As will be seen, this surely was at the heart of the Computation Center’s frontier spirit.

Recognizing that many smaller institutions lacked the capital for expensive mainframes, UT Austin acted as a regional hub through the Southwest Regional Educational Computer Network. Originally funded by the NSF, this network installed communication lines and remote terminals across Texas-scale distances, eventually expanding into a self-sustaining cooperative of twenty-three higher education institutions. To broaden the reach of these successes, Warlick helped found the Conduit coalition, a national consortium designed to evaluate, standardize, and distribute computer-related curriculum materials. Conduit solved a major technical challenge of the era by defining standards that allowed high-quality educational software written at one institution to be easily adapted and run on completely different hardware architectures at other schools. All of these efforts pointed towards 1977 when UT Austin became the first Arpanet site in Texas. This was the dawn of the Internet in Texas and in the broader region, and the story played out on the UT Austin campus.

Federal funding played an important role in the developments at UT Austin, especially via the National Science Foundation. From 1969 to 1975, two related and highly visible NSF programs broadened the user community around the CDC 6600, and then later helped introduce the DEC-10. First came the physical infrastructure of the Southwest Region Educational Computer Network, and in 1971 it was extended into the software oriented Project Computer-Based Education. 

By 1973, all of this was anchored by a centralized dual-mainframe infrastructure based on the CDC 6600 and its newer CDC 6400 sibling. This pair eventually were nicknamed the Dual Dinosaurs, and the two machines shared a common file system and mass storage, creating a robust dual-core environment where the CDC 6400 was dedicated to 128 time-sharing ports while the 6600 managed batch processing through 16 remote terminals. Roughly eighteen thousand undergraduates were utilizing these facilities, signaling that computing had moved from specialized engineering and science research into the infrastructural core of the general curriculum. This transition represented a fundamental change in university policy, as the institution began to treat computational access not as a scarce research luxury, but as a critical instructional utility. The network and the software were symbiotic arms of the same mission, to prove that computers could revolutionize teaching.

\section{RAND, SAGE, Timesharing}

The introduction of timesharing on the UT Austin CDC 6600 was a student-led initiative that fundamentally changed how the university's computers were used. Instead of originating from faculty or administration, the push for timesharing came from three graduate students. Forest Baskett was one of the instigators, and he was soon effectively running day-to-day aspects of the Computation Center systems programming staff. This was one reflection of the motto displayed in Charles Warlick’s office. Frustrated by the CDC 6600's primitive batch operating system, the students proposed replacing it with a timesharing system that utilized online terminals. They took their idea to Jim Browne, an early computer science faculty member, who supported the project and got the Computation Center's approval. 

Baskett drew inspiration from a summer job in 1964 at the System Development Corporation. Spun off from RAND Santa Monica in 1956 to handle the unprecedented software demands of the SAGE project, SDC is widely considered the world's first independent computer software company. The SDC facility originally required an air-conditioning system powerful enough to cool 20,000 homes to offset the heat generated by its early vacuum-tube systems. In this environment, operators monitored radar scopes in dimly lit rooms, surrounded by massive walls of neon bulbs displaying the state of the machine's logic gates. This summer job was while Baskett was an undergrad at Rice University in Houston and already involved with interesting computer research. He was working for a chemistry professor, running simulations of the molecules in a gas and making movies of the results. They had a cathode ray tube with a 16-millimeter film camera attached to it. It could put dots on the screen, take a picture, clear the screen, and advance the film by one click. During this period in the early sixties, Baskett was fortunate enough to experience the SAGE system at SDC.

SAGE, initially designed in the fifties as a military command-and-control system for Soviet bomber defense, featured pioneering real-time processing and early timesharing capabilities. The machine Baskett interacted with was the AN/FSQ-32, commonly referred to as the Q-32, a transistorized prototype that succeeded the massive, vacuum-tube-based machines originally built for the SAGE air-defense network. The Q-32 occupies a unique and somewhat ironic place in computing history. Its cancellation as a military asset is exactly what allowed it to become a pioneering testbed for modern interactive computing.

The Q-32 was originally commissioned by the Air Defense Command to solve the glaring vulnerability of the massive, above-ground AN/FSQ-7 SAGE blockhouses. It was designed to be installed in hardened, underground nuclear bunkers capable of withstanding 200 psi of blast overpressure, called Super Combat Centers. By 1960, RAND and the Department of Defense realized that the rapidly increasing yields of Soviet nuclear weapons and the shift toward Intercontinental Ballistic Missiles rendered even these underground bunkers vulnerable. As a result, the Super Combat Center program was cancelled, and the Q-32's military career was terminated before series production could begin, leaving the single completed prototype at the SDC headquarters in Santa Monica.

Because it was no longer needed for active-duty air defense, the Advanced Research Projects Agency, under the guidance of Joseph Licklider, repurposed the Q-32 prototype to research multi-user interactive computing. When ARPA originally funded research on the surplus Q-32, the strict goal was to study military command-and-control problems. Licklider directed researchers to curtail the specifically military projects and instead focus entirely on building a general-purpose time-sharing system. He believed that achieving Man-Computer Symbiosis, where humans and computers could interact conversationally in real time, was an absolute prerequisite before any effective command-and-control system could be built.
With the Q-32, terminal connections were routed through a frontend PDP-1. This was done to solve a major bottleneck. If the mainframe had to halt its primary arithmetic logic unit to process every single character typed on a remote teletype keyboard, its throughput would have collapsed. By placing an 18-bit DEC PDP-1 minicomputer between the mainframe and the terminals, researchers created a pioneering hierarchy of processors. The PDP-1 handled the low-speed line management, multiplexed the incoming character streams, and transferred complete words or lines to the Q-32 via high-speed parallel channels. This task consumed only 1\% to 2\% of the PDP-1's capacity, leaving the Q-32 entirely free for heavy user calculations and allowing it to successfully support 30 or more simultaneous users.

Because the Q-32 was originally designed as a single-purpose air defense machine, it lacked the hardware memory management or paging registers needed to work with multiple active programs in its 65,000-word core memory. To achieve a fluid, conversational interactivity, the timesharing monitor relied on a quantum-based round-robin scheduler. Each program was allocated a precise slice of CPU time, called a quantum. If the task wasn't finished when its time ran out, the system executed a full-image swap. It wrote the program's entire core image out to the high-speed magnetic drums, read the next program's image from the drums into the newly vacated core memory, and placed the interrupted program at the back of the queue. Researchers managed to make this heavy data swapping virtually imperceptible to the user by leveraging the parallel-bit architecture of the drum system, which featured access times of just 11.75 microseconds.

This was the specific architecture that allowed Baskett to sit at a terminal and interact with the machine. It demonstrated that a computer could be an immediate, conversational medium for creative exploration. Given the unstructured nature of his summer job, Baskett used the time to teach himself John McCarthy’s Lisp from a textbook and wrote a custom Lisp interpreter directly on the SAGE machine. This hands-on experience became his mental model for how computing should ideally operate when he later encountered the CDC 6600's restrictive, punch-card-based SCOPE batch system. The Q-32's influence extended far beyond Baskett's individual career. In October 1965, the Q-32 in Santa Monica was directly linked via a dedicated dial-up telephone line to the TX-2 computer at MIT's Lincoln Laboratory. This connection marked the first successful transcontinental exchange of data between two independent operating systems, successfully proving the viability of wide-area distributed computing and serving as a direct precursor to the Arpanet.

Understanding this context highlights exactly why Baskett found UT's CDC 6600 so frustrating just a few years later. While the Q-32 at SDC was designed for real-time command, control, and multi-user interaction, the CDC 6600 was engineered purely for maximum scalar floating-point performance. The manufacturer-supplied SCOPE operating system was strictly batch-oriented to keep the central processor constantly fed with heavy-duty numerical models and simulations. When Baskett and his fellow graduate students proposed building a timesharing system, they were essentially attempting to graft the interactive, user-friendly philosophy he had experienced at SDC onto the raw, unyielding computational power of a machine designed solely to crunch numbers. The students' frustration with the SCOPE system was entirely justified. SCOPE was designed purely for batch processing, completely isolating the user from the machine. To fix even a minor bug, researchers were forced to submit physical decks of punch cards to operators and wait hours for printed results. The system enforced counterintuitive, rigid rules, such as requiring users to define their maximum runtime in octal seconds, capping execution at exactly 77777 octal seconds, or about nine hours.

When Baskett and his two fellow graduate students approached Jim Browne with their radical idea, Browne's response was enthusiastically pragmatic: "Hmm, that could be fun. Let's try". Browne's backing was the critical catalyst for the project. He had to navigate the university's administrative hierarchy to convince the Computation Center to allow a small group of students to completely replace the core software of a \$5.9 million supercomputer. In addition to his work in systems software, Browne served as a Principal Investigator for the Conduit project at UT Austin, working alongside Charles Warlick and George Culp to test, evaluate, and distribute computer-based curriculum materials across different universities. Ultimately, his early experiences supporting timesharing and multi-institutional resource sharing shaped his later career, and Browne went on to become a major proponent of national high-speed computer networks.

Once Browne secured the Computation Center's approval, Baskett and his team ingeniously repurposed the CDC 6600's unique hardware to solve the software bottlenecks. By programming the mainframe's ten independent Peripheral Processors to handle the input and output operations of remote interactive terminals, they freed the central processor to execute user programs in rapid, multiplexed time slices. The trust that Jim Browne, Charles Warlick, and the Computation Center placed in these students yielded extraordinary results. The initial system was up and running within a year and a half, and its subsequent revisions proved so stable that it remained in active production at the university for a remarkable ten years. Browne also went on to serve as Baskett's doctoral thesis advisor as part of supporting the research on system scheduling and queuing theory that emerged from the project.

The primary goal of the new UT system was to make the computer easier to use, more enjoyable, and highly productive for its target audience of faculty researchers and graduate students. Through data analysis of user habits, Baskett deliberately aimed to optimize the system to minimize customer complaints. The system initially used teletypes and relied on the existing compilers supplied by CDC, ensuring the operating system maintained all the interfaces that users were already accustomed to. Baskett implemented a job-scheduling method using round-robin timeslicing, similar to the one he had experienced on the SDC Q-32, keeping the time slices as small as possible while remaining consistent with system overhead. 

Tasks were kept memory-resident and managed via the 6600's base and bounds registers, which provided memory protection on a per-job basis. Because the CDC 6600 lacked hardware paging, segmentation, or virtual memory mapping, user programs were forced to reside in contiguous physical blocks within the central memory. The base and bounds registers provided strict per-job memory protection so users couldn't maliciously or accidentally corrupt each other's data. To effectively multiplex dozens of users with round-robin timeslicing, the system had to swap these memory blocks rapidly. The team achieved this by leveraging the extended core storage system. When a user's time slice expired, the TAURUS scheduler initiated a fast block transfer, copying the user's entire contiguous address space into extended core and immediately swapping the next active user's program into central memory.

During this period, Baskett encountered Seymour Cray and asked the legendary architect about modifying the 6600 hardware so that privileged instructions in user mode would cause an exception rather than a no-op. This was essentially a plea for hardware-level virtualization. Cray’s succinct refusal "No, I don’t think so" illustrated the persistent gap between architectural vision and hardware implementation that Baskett would spend his career bridging. As he assumed leadership of the Computation Center’s twenty-five person systems programming staff, he balanced these practical infrastructure challenges with the theoretical rigor that would define his doctoral dissertation. Between 1971 and 1982, his career exemplified a unique industrial-academic synthesis, bridging the gap between national laboratories and corporate research. He led the Demos operating system for the Cray-1 at Los Alamos National Laboratory, which was notable for its use of software-based property tags, an early precursor to modern object-oriented systems. Simultaneously, he conducted VLSI research at Xerox PARC.

The initial timesharing service, known as RESPOND, was officially initiated on the 6600 in March 1967. It proved to be a major success, with the system and its subsequent revisions remaining in production for a decade. RESPOND was later replaced by a more advanced system called TAURUS (Texas Anthropocentric Ubiquitous Responsive User System), which operated as an integral part of the UT-2D dual operating system, managing both the CDC 6600 and 6400. The UT Austin project was highly influential in the broader computing industry. It demonstrated the viability of timesharing on the CDC 6600, prompting both Control Data Corporation and the Lawrence Livermore National Laboratory to realize the need for such systems and launch their own competing efforts.

\section{Southwest Network}

In the sixties, computer timesharing and networking, packet switching in particular, evolved quickly as associated areas of research. Computing resources were scarce and expensive, and the overall motivation was to share the available resources among more users, with both timesharing and packet switching slicing the resources into increasingly granular pieces for better distribution both temporally and spatially. Timesharing was the immediate, urgent problem that needed solving, and packet switching was born partly as a tool to expand and scale timesharing upward. All of these developments were bound up within the stories of the SAGE Q-32 in Santa Monica during the mid sixties and then the UT Austin CDC 6600. Essentially, the evolution of timesharing on the Q-32 led to its central role in early computer networking and packet switching research, and was reflected and amplified in the UT Austin Computation Center.

Timesharing revolutionized computing by slicing up a single machine’s processing power, rotating between multiple users so quickly that every user felt like they had the machine to themselves. To use a timeshared computer, you needed a way to talk to it. This drove the immediate need for networking. At first, this meant connecting local terminals, usually teletypes, to the central mainframe via direct serial cables within the same building. Soon, users wanted to access these mainframes from other buildings, universities, or research labs. This forced the use of telephone lines and early acoustic-coupler modems to transmit data over long distances, but traditional phone networks used circuit switching, holding open a dedicated line for the length of an entire conversation. This was inefficient for computers because human interaction with a computer is inherently bursty. A user will type a few keystrokes and then pause for seconds or minutes to think or read output. Before timesharing, batch processing dedicated an entire multi-million-dollar mainframe to one task at a time, meaning the central processor was completely wasted if it had to wait for human input or setups. Before packet switching, the telephone network relied on circuit switching, which held a physical pathway open exclusively for the duration of a session. During a user's pauses, the expensive communication line sat completely idle, wasting a significant portion of the bandwidth. Researchers developed packet switching partly to address this inefficiency. Data was broken into small packets, sent across the network dynamically, and reassembled at the destination. This allowed the network itself to be timeshared. Multiple users could send bursts of data over the same physical wires simultaneously without clogging the system. 

The concepts of timesharing and packet switching share a profound functional and historical similarity. Both technologies were developed to take advantage of the bursty nature of human-computer interaction by dynamically slicing up an expensive resource and sharing it among multiple users. Both concepts abandoned the idea of dedicating an entire resource to a single purpose, instead breaking the resource down into manageable pieces. Timesharing slices the computer's processor cycles into millisecond intervals or quanta. Packet switching chops a user's digital message into small, fixed-size data packets, such as 1024 bits. Rather than pre-allocating a chunk of time or a dedicated wire, both systems use multiplexing to dynamically assign resources only when they are actively needed. The computer operating system's monitor rapidly cycles through active users. If one user is pausing to think, the system immediately skips them and grants processing time to the next active user's task. Meanwhile, data packets from dozens of active timesharing users are interleaved and sent over the same physical wire. If one user pauses, the network channel is not locked or wasted. It instantly routes packets belonging to other active users.

Because this slicing and switching occurs at lightning speed, both systems successfully deceive the human user. Timesharing switches between tasks so rapidly that it provides each user with the illusion of having the entire computer completely at their disposal. Packet switching performs a similar feat over the network, providing the seamless feel of a dedicated, continuous circuit without actually requiring one. Historically, these analogies were not a coincidence. Donald Davies, the co-inventor of packet switching along with Paul Baran of RAND, explicitly designed the concept in 1965 because he recognized that traditional circuit-switched telephone networks were fundamentally incompatible with the bursty traffic generated by the new interactive timesharing systems. Packet switching was directly modeled to bring the same economic efficiency and dynamic sharing to telecommunications that timesharing had just brought to mainframes.

The Q-32 was used for early long-range networking experiments. In late 1963, researchers established a 300-mile link between the Q-32’s PDP-1 front-end in Santa Monica and a CDC 160A minicomputer at the Stanford Research Institute. This connection utilized two full-duplex telephone lines. The defining feature of this experiment was the strict separation of control and data. One telephone line was dedicated to sending standard executive commands while the second line was dedicated to raw data. The data line bypassed the Q-32's executive command parser and routed directly to the active program. This setup allowed remote SRI users to interactively perform text searches on bibliographic databases stored on the Q-32's high-speed magnetic drums.

By 1969, the early experiments with the Q-32 had led onwards to UT Austin and the NSF establishing a large-scale regional network, the Southwest Region Educational Computer Network. Its goal was to provide computing power to smaller colleges and universities across Texas. This regional network was an experiment funded largely by the NSF to mitigate the prohibitive upfront costs of mainframe computing. With a simple hub-and-spoke model, the Southwest Network provided the computational plumbing necessary for remote batch processing and timesharing and created an environment for computer-based education to flourish. It functioned as infrastructure, the hardware backbone and telecommunications, that linked the central processing power and software of the UT Computation Center to a distributed web of smaller colleges, junior colleges, and secondary schools. 

The Southwest Network was centered on the Computation Center CDC 6600 and 6400. Schools such as Southwest Texas State in San Marcos, Rice in Houston, Austin High and McCallum High, and various junior colleges connected to this central hub via remote terminals, usually teletypes. It expanded from nine initial institutions in 1969 to twenty-three in 1972, with UT Austin serving as the central host. It democratized access to high-level processing power for both administrative tasks and, more importantly, classroom instruction. While the network was a technical success and demonstrated that communication lines and terminals could be stabilized across a distributed institutional landscape, it reached a definitive structural limit. The project successfully navigated the connectivity bottleneck but failed to address the more elusive problem of educational effectiveness. The network proved that data could be transported, but it lacked the content necessary to justify the technology’s presence in the classroom. This highlighted a critical gap. Technical availability did not equate to educational utility, necessitating a shift toward the production of high-quality instructional materials.
Student oriented educational software became the focus from 1971, championed by faculty such as John Allen, Joseph Lagowski, and George Culp. The efforts centered in the Project Computer Based Education office in the aerospace building, still referred to as the Engineering Sciences Building at the time but changed to WRW soon after. The project was a high-level initiative at UT Austin focused on education and researching how students learn through computer-assisted instruction. It was designed as a rigorous five-year curriculum development effort, 1971 through 1975, involving seventy-five professors and over 4,000 students in the creation of actual classroom modules. 

The project was also conceived as a targeted intervention designed to infuse NSF capital directly into the production of instructional material to stimulate a market and bypass the traditional textbook-publishing model, which had historically suffered from institutional inertia and a failure to accommodate the complexities of software. By funding the production of quality modules, the project sought to provide the empirical evidence of effectiveness required to jumpstart the entire educational technology ecosystem. It engaged the Human Resources Research Organization for data collection and independent evaluation. This was also part of a democratizing effort, pushing computer usage into the soft sciences, including linguistics, psychology, and home economics. This disciplinary expansion validated the project’s educational theories and successfully transformed university policy by embedding instructional computing into the humanities and social sciences.
 
Probably the most significant artifact of the instructional ecosystem was the Symbolic Interactive Matrix Processing Language, or Simple, a direct ancestor of the commercial Matlab environment still common in engineering today. Conceived as a conversational, matrix-oriented programming language, Simple was designed as a bridge for beginners in engineering who needed to work with sophisticated matrix calculations despite minimal prior computational experience. The focus was on computer-based teaching techniques in undergraduate science and engineering education and transforming the teaching of linear algebra and numerical analysis. It was generally used via remote terminals, specifically acoustically coupled Datapoint terminals and standard teletypes, which allowed students to bypass the delays of traditional batch computing. It included a Teach command that was central to the instructional shift, providing self-paced, inquiry-based documentation that enabled students to obtain quick results and proceed at their own rate. For the mature user in aerospace or structural engineering, Simple provided rigorous tools for complex operations, including essential matrix decompositions. By accommodating both batch and timesharing modes, Simple ensured its own survival and grounded high-level engineering research in a versatile, interactive everyday tool. 

The workaday success of Simple and Matlab has demonstrated the real-world value of this type of code. Despite a distinct lack of glamour and hype, they endure because they reflect a deep connection between engineering and computational modeling and simulation. This special corner of the software world and its particular links to the engineering world are well worth further exploration. A quote from Matlab’s creator Cleve Moler captures this, here he's discussing the course he taught at Stanford circa 1980.

\begin{quotation}
The computer science students and numerical students were not very impressed. This was not a sophisticated language, they were taking courses from Don Knuth and John McCarthy, and so on, who did real computer science. There was not much mathematics going on here, these algorithms were old, traditional algorithms, there was no convergence theory, no partial differential equations, this was not the usual fare for the graduate course in matrix computation. The other half of the students in that class were students that had come from various departments in engineering, electrical engineering, mechanical engineering, nuclear engineering, they knew about matrices, they’d used matrices a lot in their theoretical work, many of their textbooks described things in terms of matrices. They had done computation with matrices using Fortran. Some of them knew about LINPACK and EISPACK, when they saw MATLAB they thought this was terrific.
\end{quotation}

Moler created Matlab in the mid seventies at the University of New Mexico in Albuquerque and it would be good to learn how much influence Simple had on the beginnings of Matlab. There were definite connections between UT and New Mexico, particularly in the context of White Sands and Los Alamos.

When researchers at UT Austin developed instructional software such as Simple or an interactive chemistry simulation, they could deploy the modules to a community college hundreds of miles away in real-time. It addressed the what of instruction, developing sophisticated computer-assisted instruction modules that moved beyond simple rote memorization into the realm of interactive problem-solving. This was particularly evident in the hard sciences and mathematics, where project researchers utilized the CDC 6600’s capabilities to create simulations and drill-and-practice routines that were previously impossible on standalone systems. 

The Southwest Network allowed for the dissemination of standardized curricula across diverse geographic and socioeconomic boundaries, effectively proving that a centralized mainframe could serve as a universal tutor and making UT Austin a central node of educational innovation and outreach to the community. The partnership of networking and software proved that computer-assisted instruction was viable on a large scale. It helped standardize educational software across the region and on a national scale via the Conduit program, extending to Dartmouth, the University of Iowa, North Carolina, and Oregon State, and demonstrated that the high cost of mainframe computing could be offset by sharing the resource across multiple educational institutions. And it solidified UT Austin’s reputation as a national leader in academic computing, a precursor to Austin's current status as a tech mecca. 

The historical trajectory of academic computing in the early seventies represented a shift from autonomous one-campus one-computer models toward a centralized, shared resource paradigm. This evolution was not just a matter of technological convenience but was driven by institutional priorities as statewide governing boards and coordinating councils sought to impose order on burgeoning computational demands. As analyzed in Charles Mosmann’s 1974 study on academic computer planning, this era was defined by a network bandwagon phenomenon, where planners pursued the perceived economic efficiencies of large-scale star networks. The Mosmann critique highlighted a fundamental fallacy in this high-level top-down planning. Merely installing communication lines to an existing computing center did not constitute a functional network. Without stability, application quality, and localized consulting services, centralized systems often suffered from chronic under-utilization, frequently seeing less than five percent off-campus use. 

While the early years of the seventies opened up new technological horizons, by 1974 many researchers experienced a closing in of options as autonomy was sacrificed for perceived cost effectiveness. These tradeoffs pointed the way towards future decentralized networks and user communities. The proliferation of video terminals from the mid seventies onward also rapidly changed the user experience, providing much higher levels of interactivity than the teletypes common in the early seventies, and opening entirely new creative channels, especially as graphics terminals became affordable. All signs were in fact pointing forward towards the Internet, which was already growing rapidly in its preliminary form as the Arpanet, and was soon to arrive at UT Austin in 1977, the first node in Texas and the wider region.

\section{ARPANET And Internet}

Much as it had done for timesharing and interactive real-time computing, the Q-32 in Santa Monica foreshadowed the 1977 dawning of the Arpanet and Internet at UT Austin and in Texas. The Q-32 was retired around 1970 and did not itself serve as a node in the operational Arpanet, but it played a critical role as a direct technical ancestor. In 1965, ARPA funded a project proposed by Thomas Marill and overseen by Larry Roberts to link the Q-32 directly to the TX-2 computer at MIT's Lincoln Laboratory in Massachusetts. This experiment proved that two independent, geographically distant timesharing operating systems could exchange digital data and invoke programs remotely. It also exposed severe flaws in the era's networking capabilities. These technical frustrations directly convinced Roberts, who would subsequently become the program manager and principal architect of the Arpanet, that building a robust, large-scale computer network would require abandoning circuit-switching in favor of packet-switching technology.

Marill and Roberts devised what they called the "elementary approach". The primary goal was to bridge the gap between the two structurally incompatible operating systems, the TX-2's APEX and the Q-32's TSS, without having to modify or rewrite their complex core kernels. Under this protocol, the local user executed an application that intercepted terminal inputs, repackaged them, and directed them over the communications link so that the remote host monitor treated the connection exactly as if it were a local user terminal. A successful demonstration of this setup involved a researcher at the TX-2 in Massachusetts utilizing a program called the Algebraic Translator to automatically dial the Q-32 in California. The program bypassed standard human login prompts to gain administrative access, loaded a Lisp compiler on the Q-32, transmitted a complex Lisp program across the country, and had the Q-32 execute the calculations before returning the results to the local TX-2 console.

The experiment was a technical success but a practical headache, proving two major things. A user on the TX-2 in Massachusetts could log into the Q-32 in California and run a time-shared program. It showed that computers didn't just have to talk to human typists, they could talk directly to each other and share workloads. And it highlighted the fatal flaw of using standard telephone infrastructure for computing. Because the connection used circuit switching, with a dedicated analog line, it was unreliable and inefficient. Because human-computer interactions and data exchanges happen in short bursts, the line remained idle for most of the session, resulting in extremely low bandwidth utilization and proving that circuit-switched computer networks would be prohibitively expensive. Furthermore, the analog lines were highly susceptible to electrical noise and signal attenuation, which frequently caused data corruption through bit-flipping. Since the remote software did not have automated, system-level error detection and recovery, a single flipped bit meant the user’s local software had to abort the session, clear remote memory buffers, and retransmit the entire data block.

Bob Taylor, the director of ARPA's Information Processing Techniques Office and a UT Austin graduate, recruited Roberts from Lincoln Laboratory to become the program manager and principal architect of the Arpanet in Washington. Remembering the lessons of the 1965 Q-32 and TX-2 hookup, Roberts realized that instead of dedicated phone circuits, the new network had to use packet switching to handle the data efficiently. Some researchers advocated for the dual-line model used in the 1963 Q-32 and SRI experiment, arguing that separating command and data channels simplified hardware and eliminated processing overhead. Roberts argued that leasing dual transcontinental telephone lines was economically unsustainable for large-scale networks. He insisted that any viable network must use a single physical channel, with the operating system dynamically multiplexing and parsing commands and data. The 1965 experiment's severe telecommunications bottlenecks had decisively proven that scaling a wide-area network required abandoning circuit-switched telephone lines in favor of packet-switching technology. By breaking data into small, self-contained packets, the network could maximize line utilization and automatically handle error recovery through intermediate switches, rather than forcing the user's application program to handle all error checking.

Drawing on his experience connecting the TX-2 and Q-32, Roberts initially proposed at a 1967 meeting that all host mainframes connect directly to one another and manage their own network administration, terminal routing, and error checking. Mainframe operators, termed principal investigators in the ARPA context, fiercely opposed this plan. They were highly protective of their expensive mainframes and refused to sacrifice 10\% to 15\% of their limited computing power to handle network overhead. 

In response to this resistance, computer scientist Wesley Clark told Roberts, "You've got the network inside out," and proposed a decoupled architecture. Clark suggested installing a small, standardized minicomputer at each site to act as a universal interface. These minicomputers would all speak a uniform language and handle all the routing, packet buffering, and error-checking tasks, completely insulating the main host computers from the network overhead. Roberts adopted the idea and named the specialized communications processors Interface Message Processors or IMPs. By separating communication logic from application processing, the IMPs allowed structurally incompatible mainframes to easily join a unified wide-area network, establishing the foundational architecture of the Arpanet and the modern Internet.

The IMPs became the defining feature of the Arpanet by solving the critical problem of mainframe incompatibility and network overhead. Instead of forcing heterogeneous mainframes to manage the network, ARPA contracted Bolt, Beranek and Newman to build small, standardized minicomputers, the IMPs, to act as universal interfaces. Built from ruggedized Honeywell DDP-516 computers, they were the world's first true network routers. They insulated the host mainframes by taking over four foundational tasks: accepting raw data from the connected host computers, breaking the data into small 128-byte packets for transmission, attaching sender and receiver address labels to each packet so they reached the correct destination, and using continuously updated routing tables to send packets through the fastest, most efficient paths based on current network traffic. By offloading error checking, routing, and packet buffering entirely to the IMPs, the network allowed structurally incompatible mainframes to easily plug into a unified wide-area network via a single, simple digital interface.

In the early days of the Arpanet, researchers recognized that the network needed to allow users to log into remote systems, and enable the movement of files between different machines. To meet these needs, Telnet and the File Transfer Protocol were developed as the Arpanet's first major host-to-host application protocols. Once the foundational Network Control Program was implemented in 1970 to handle reliable, bidirectional communication between host computers, higher-level applications like Telnet and FTP could be built on top of it. These protocols operated at the application layer of the Arpanet and functioned independently of the underlying network service, serving as an early example of protocol layering.

Telnet was invented by the Network Working Group in 1969 and outlined in RFC 15, providing a virtual teletype connection to a remote host. This allowed researchers to remotely log into computer systems to access timesharing interfaces. FTP was specified in 1971 with RFC 114 and fully implemented by 1973 with RFC 354. It enabled users to access, copy, save, and transfer data files across the network. In addition to handling standard file transfers, FTP was also initially used to enable the Arpanet's rudimentary electronic mail. By providing standardized methods for remote login and file transfer, these protocols proved that connecting diverse computers over a wide-area network was far more economical than purchasing and duplicating expensive mainframe computers at every research site.

While the IMPs provided the physical and logical infrastructure, the network's cultural breakthrough occurred at the Arpanet Ball. The Ball is the nickname for the International Conference on Computer Communications held in October 1972 at the Washington D.C. Hilton. Organized by Bob Kahn of BBN, this event was the first large-scale public demonstration of the Arpanet. Kahn's team set up a node with 40 terminals connected to a Terminal IMP, allowing attendees to interact directly with the network. Attendees were handed a booklet containing scenarios. These were step-by-step instructions that allowed them to remotely log in to different computers across the country and test distributed applications themselves. One popular demonstration was an early computer-to-computer chat between two AI programs: a simulated psychotic patient named PARRY at Stanford discussing its problems with a simulated therapist known as the Doctor at BBN.

The Arpanet Ball successfully proved to the international research community that packet-switching and remote resource sharing actually worked, and triggered new growth in the network's size and enthusiasm in 1973. It was recognized that there was an immediate need to establish agreed-upon protocols for interconnecting independent networks globally. As a direct result of a meeting at the 1972 ICCC, the International Network Working Group was formed with Vint Cerf as its first chairman. The INWG spearheaded the collaborative effort to design an open-architecture network, paving the way for Cerf and Kahn to develop the TCP/IP protocols that would ultimately transform the Arpanet into the modern Internet.

The first IMP in Texas was delivered to UT Austin in 1977 and initially installed alongside the DEC-10 in the HRC. This was the dawn of the Internet in Texas. Because the standard Tops-10 operating system did not yet support interfacing with the IMP, the DEC-10 did not become the first Internet computer in Texas. The many PDP-10s across the Arpanet from its earliest days in 1969 were running custom operating systems, and the HRC machine was kept factory stock. Instead, the first actual host to connect was a PDP-11/45 running a hacked version of UNIX, located across campus in the Daily Texan composing room. Accessible via dial-up modems by a small group of users, this machine became the UTEXAS host on the net and served as UT's pioneering gateway. Here is a full quote from Clive Dawson, who was involved in all of these events, particularly for the DEC-10 and DEC-20.

The IMP arrived on the UT campus in 1977 and was installed alongside the DEC-10 in the HRC. This was because originally there were hopes that the DEC-10 would become the first Arpanet host on the Texas node. But DEC never offered support for the Arpanet on the DEC-10’s standard OS, Tops-10. The only DEC-10 OS's that supported a connection were TENEX (developed at BBN), ITS (the Incompatible Timesharing System at MIT), and several heavily customized versions of Tops-10 at places like CMU, Harvard, and Wharton. The Stanford AI Lab ran WAITS (the West-coast Alternative to ITS) on their DEC-10s. So the DEC-10 in HRC was never connected to the Arpanet. Meanwhile, there was a PDP-11 running UNIX in the Daily Texan composing room at TSP (Texas Student Publications). A hacked version of UNIX from the U. of Illinois was installed, and the IMP was moved across campus to the Communications Building. Thus, a PDP-11/45 running UNIX became UTEXAS, the first Arpanet host on campus. The select few who knew about this would connect to the 11 via dialup modems and from there could venture out to the rest of the Arpanet.

Sometime around 1978-80, a newer IMP was delivered by BBN and installed in the Computation Center. A PDP-11/70 was installed there to connect to the new IMP, and so the UTEXAS host was rehomed from TSP to the Comp. Center. It was only when the first DEC-20 arrived and was installed in Painter Hall that the PDP-11 became UTEXAS-11 and the 20 was named UTEXAS-20. As I recall, the 11 retained the host nickname of UTEXAS.

At about the same time, specialized research hardware was acquired by the Certifiable Minicomputer Project (CMP) for the purpose of providing secure encrypted communications. This required that a very thick cable with 32 twisted pairs of copper be run from the 21st floor of the tower, down the elevator shaft and through the steam tunnels to the Comp. Center, where it plugged into the IMP. When the CMP project was completed several years later and the hardware decommissioned, the cable was left in the elevator shaft, as it was not worth the trouble to remove it. I wonder if the folks currently remodeling the Tower have come across it?! (Hat tip to Rich Cohen and Clyde Hoover for filling in some details of this story.)

Before the DEC-10 was installed, the UT Austin Computation Center remained dominated by the legacy of the CDC 6600, a machine designed for high-throughput floating-point computational science. The CDC 6600 operated under UT-2D, a specialized operating system developed locally by the Computation Center staff to maximize the efficiency of its 60-bit word architecture. There were locally developed timesharing extensions for the 6600, but at their heart, CDC and Cray machines were intended for the types of dedicated floating-pointing linear algebra based modeling and simulation common in the hard sciences and engineering. Likewise, the Fortran programming language was intended and designed specifically for these types of problems. Floating-point linear algebra based problems were intentionally an ideal match with the Fortran language and the CDC and Cray hardware, perhaps with a small amount of assembly language thrown in for optimization within deep inner loops. 

This was the legacy and environment in which the DEC-10 arrived. The coexistence of these systems allowed the university to partition its workload effectively. The CDC machines remained the workhorses for massive simulations, while the DEC-10 became the preferred platform for software development, interactive modeling, and creative or social computing. The introduction of the DEC-10 at UT Austin in 1975 represented a fundamental shift in the university's approach to centralized computing, transitioning forward from the more rigid, batch-oriented legacy of the sixties and toward a more fluid, interactive paradigm. 

\section{HRC And The DEC-10}

In 1975 the university embraced the DEC-10's 36-bit architecture and prioritized timesharing, user accessibility, and real-time interaction, shifting the campus computing culture toward a more exploratory paradigm. At UT, the DEC-10 era brought together a diverse cast of users, ranging from creative coders who reveled in the complex, messy reality of systems programming, to algorithmic formalists who demanded provable logic. A range of subcultures found a suitable home and shared laboratory on the DEC-10 and its Tops-10 timesharing operating system. This machine created an environment that managed to simultaneously support a new diversity of projects and people. It seems appropriate that the DEC-10 was located in the Harry Ransom Center, an arts and humanities landmark, rather than the traditional old Computation Center machine room with its CDC hardware.

A good indicator of the spirit of the HRC location was that the offices housing the DEC-10 and staff in the HRC were adjacent to offices of the Linguistics Research Center. Another quote from Clive Dawson gives a vivid sense of the HRC environment.

\begin{quotation}
Rich, the person you are thinking about across the hall at the Linguistics Research Center (LRC) was Dr. Helen Jo Hewitt, aka “HJ”.  She was a BIG fan of TECO, and was one of my “power” users.  She was one of the early “font hackers”, and I had developed a bunch of TECO macros for her.  She would come across the hall and pose an interesting problem, and a few minutes, or hours, or days(!) later, I would present her with another TECO macro for her growing collection. These tools would enable her, for example, to use a meta-language of her own creation to insert diacritics into a document, and then apply the appropriate macro to convert the meta-language into actual escape-sequence commands for the robotic typewriter to backspace, make micro-adjustments to the carriage, and type the appropriate symbol(s) to produce the desired diacritic at j-u-s-t the right position. Later on, as her daisy-wheel collection of different type-fonts grew, a macro would be used to type a partially-filled page of text, then command the typewriter to roll the page back to the beginning and pause. She could then switch in, say, the Greek daisy-wheel, after which the job would continue with the typewriter skipping directly to each blank spot where a Greek character needed to be typed.
\end{quotation}

From an administrative perspective, the DEC-10 offered an attractive alternative to the more expensive mainframes of the era. The system was designed to be compact and efficient, requiring significantly less site preparation than the CDC 6600 or comparable IBM 360 systems. The upfront purchase price was less than a million dollars, which meant a lower capital outlay than traditional mainframes. Site prep cost and specialized construction were relatively minimal and estimated at less than \$20,000. Floor space, cooling, and staffing needs were also relatively low. Notably, one operator per shift was still expected, and someone would be with the machine in HRC around the clock. The efficiency of the Tops-10 operating system and its reentrant compilers allowed a single copy of a compiler such as Fortran to serve many users concurrently, conserving central memory and reducing the need for constant swapping to disk. This architectural efficiency was a primary reason why the university could support a “whole campus of users" on a system that cost less than half as much as its competitors.

For Don Good, Woody Bledsoe, Jim Browne, and the Institute for Computing Science and Computer Applications, the DEC-10 offered a large address space and stable multitasking capabilities. Their work aimed to treat computer programs as formal mathematical objects, completely replacing empirical, trial-and-error debugging with rigorous mathematical proofs. ICSCA was located on the 20th and 21st floors of the UT Tower. In the seventies and early eighties, the major research thrust of ICSCA was program verification. Motivated by deep concerns over the unreliability of software controlling critical infrastructure and nuclear armaments, the group developed the Gypsy Verification Environment, an interactive system running on the DEC-10 that transformed programs and specifications into logical formulas to be mechanically proven. 

In the Certifiable Minicomputer Project of the late seventies, the group used the DEC-10 and Gypsy to achieve a milestone for the era by mathematically verifying Arpanet security code associated with historically important encryption hardware from Bolt Beranek and Newman. Developed by BBN starting in 1973, the Private Line Interface aimed to give Arpanet users the equivalent of a private, leased line over a shared public infrastructure. The PLI pioneered the concept of edge cryptography for packet-switched networks, establishing an architectural trajectory that eventually evolved into today's IPsec and High Assurance Internet Protocol Encryptor standards. 

To operate at the edge of the network without requiring extensive modifications to the existing Arpanet infrastructure, the PLI relied on a breakthrough known as selective payload encryption. By encrypting only the sensitive message contents while leaving the machine-readable routing and addressing headers in plaintext, the PLI allowed classified traffic to be securely routed through shared IMPs without requiring the internal network switches to hold any cryptographic keys. The devices were approved by the NSA in 1975 for limited deployment to protect classified data and originally utilized manually keyed cryptographic units. The success of the PLI's edge-based architecture quickly spawned a lineage of advanced successors that shaped the modern internet. 

The physical, single-purpose black box edge devices engineered by BBN inspired UT Austin's Certifiable Minicomputer Project. Because dedicated hardware like the PLI was rigid, defense agencies wanted to know if a more flexible, software-defined gateway running on a minicomputer could provide the exact same level of edge security. By successfully writing and mathematically verifying a 5,000-line cryptographic gateway in the Gypsy programming language, the UT Austin team proved that software could sit at the edge of the network and securely interoperate with BBN's physical hardware without leaking data, proving that absolute security guarantees could be enforced in software at the network's edge. As described in the quote from Clive Dawson, a PLI was installed and required that a very thick cable with 32 twisted pairs of copper be run from the 21st floor of the tower, down the elevator shaft and through the steam tunnels to the Computation Center, where it plugged into the IMP. When the CMP project was completed several years later and the hardware decommissioned, the cable was left in the elevator shaft, as it was not worth the trouble to remove it.

The work on the DEC-10 did not occur in isolation. A critical component of the university's software library came from the Digital Equipment Computer Users Society (DECUS), a global organization that facilitated the exchange of programs and technical insights. The DECUS library tapes provided the UT Austin community with hundreds of specialized programs in both assembly and Fortran. The software exchange service provided by DECUS allowed UT Austin's systems staff to communicate directly with their counterparts at other major installations. This exchange of suggestions, hints, bugs, fixes, and program plans was vital for maintaining the stability of the local DEC-10 environment and for keeping the university's software offerings at the cutting edge of the technology.

The success of the DEC-10 at UT Austin was not just a result of the software but also of the operational management provided by the Computation Center. The center’s staff had to balance the needs of a diverse user base, ranging from students learning Basic to doctoral candidates running complex Fortran simulations. From an administrative perspective, the DEC-10 offered an attractive alternative to the more expensive mainframes of the era. The Computation Center implemented policies that fostered a sense of community around the DEC-10. The system's Bboard and GRIPE utilities became vital communication channels, allowing users to interact with the system staff and with each other. 

The staff itself was described as an extremely close-knit group that felt a deep personal connection to the machine. This connection was shared by the users, many of whom spent thousands of hours connected while developing software and helping others. This social aspect of the DEC-10 was a direct precursor to the modern online communities that would emerge with the growth of the Internet, a network that the DEC-10 itself helped to develop. Despite its popularity and technical excellence, the DEC-10 at UT Austin eventually faced the inevitable reality of technological obsolescence. By the early 1980s, the computing world was moving toward 32-bit virtual memory architectures, most notably represented by DEC's own VAX series. 

The decommissioning of the UT Austin DEC-10 on October 31, 1982, was a momentous event that marked the end of an era for the university's computing community. The decision was driven by the chronic state of saturation common to university computer systems, combined with the emergence of newer machines like the DEC-20, which offered better integration with the Arpanet. The final night of the DEC-10 was documented by Clive Dawson in his piece "The Soul of an Old Machine". The account describes a scene of both celebration and mourning, as staff members and users gathered to witness the final hours of the system. Users sent final messages of gratitude to the machine, reflecting the profound impact the system had on their academic and personal lives.
